\chapter{Phase 2: decomposition and regularity}\label{ch:cII2}
\chapterbrief
{Part~I, Chapter~\ref{ch:cI2} (A10--A18): the finite quotient, the twisted Poincar\'e lemma, the factor decomposition, the regularity criterion and density.}
{to prove the remaining decomposition results: the one-dimensional lemmas on rectangles, constructive projectors, the Vandermonde formula, interior regularity, the cases where the $H^m$ split holds, the cohomological criterion for regions with holes, and the closedness results for families of transports. The chapter ends with the complements on variable-coefficient identification proved in the same research package.}
{Dim.~B22--B42.}

%=====================================================================
\section{Rectangles and convex sets}

\begin{dimcard}{B22}{C12-L1}{One-dimensional lemmas on a rectangle}{\proved}{}
On $\Omega=I\times J$: (a) $\{\varphi\in\mathcal D'(\Omega):\partial_x\varphi=0\}=\{c\otimes1:\ c\in\mathcal D'(I)\}$; (b) there is a linear $\mathcal J$ with $\partial_x\mathcal J=\mathrm{id}$ and $[\partial_t,\mathcal J]=0$; (c) the twisted Poincar\'e lemma holds on $\Omega$; (d) $\{w:\ D_tw=D_xw=0\}=\C\,e_{(p,q)}$.
\end{dimcard}
\begin{proof}
\step{1}{(a), (b)}
Fix $\theta\in C_c^\infty(J)$ with $\int\theta=1$. Every test function splits as $\psi=\partial_xK\psi+\Psi\otimes\theta$ with $\Psi(t)=\int\psi(t,x)\dd x$ and $K\psi$ compactly supported. If $\partial_x\varphi=0$, then $\langle\varphi,\psi\rangle=\langle c,\Psi\rangle$ with $\langle c,\chi\rangle=\langle\varphi,\chi\otimes\theta\rangle$. Define $\langle\mathcal J\varphi,\psi\rangle=-\langle\varphi,K\psi\rangle$.
\step{2}{(c), (d)}
Remove the twist by $e=e_{(p,q)}$. With $w_0=\mathcal J\tilde\psi$, $\partial_x(\partial_tw_0-\tilde\varphi)=0$, so $\partial_tw_0-\tilde\varphi=c\otimes1$ by (a); subtract a primitive $C\otimes1$ of $c\otimes1$. For (d), $\tilde w=c\otimes1$ with $c'=0$.
\whyworks{on a rectangle every equation in one variable can be integrated with the other variable as a parameter.}
\end{proof}
\uses{none.}

%---------------------------------------------------------------------
\begin{dimcard}{B23}{C12-L2}{Finite-dimensional modules on a rectangle}{\proved}{}
For a finite-dimensional $\ring$-module $M$ and $\mathcal A$ on a rectangle: $\mathrm{Ext}^1_\ring(M,\mathcal A)=0$, $\dimc\mathrm{Sol}_{\mathcal A}(M)=\dimc M$, and every solution is an exponential-polynomial.
\end{dimcard}
\begin{proof}
\step{1}{Simple modules}
For $M=\ring/\mathfrak m_z$ the Koszul resolution reduces both claims to B22(c), (d).
\step{2}{Induction}
A common eigenvector of $\tau,\xi$ on $M$ gives a simple submodule $S$; the long exact sequence for $0\to S\to M\to M/S\to0$ and induction give the two claims.
\step{3}{Exponential-polynomials}
The minimal polynomials $s(\tau)$, $r(\xi)$ of $\tau,\xi$ on $M$ kill every solution value, and A2, Steps~5--7, applies.
\whyworks{finite modules are built from points, and each point contributes one exponential.}
\end{proof}
\uses{B22, A2.}

%---------------------------------------------------------------------
\begin{dimcard}{B24}{C12-L5}{Functions constant along a direction on a convex set}{\proved}{}
If $D\subset\R^2_{(s,r)}$ is open and convex and $\varphi\in\mathcal D'(D)$ has $\partial_r\varphi=0$, then $\varphi=F(s)$ for a unique $F\in\mathcal D'(\mathrm{pr}_sD)$; if $\varphi\in L^2_{loc}$, so is $F$.
\end{dimcard}
\begin{proof}
\step{1}{Local form and gluing}
On each rectangle $Q\subset D$, B22(a) gives $\varphi=F_Q\otimes1$. Two rectangles whose projections overlap are joined by a thin rectangle along a segment inside $D$ (convexity), so the $F_Q$ agree and glue.
\step{2}{$L^2$}
$F=\int\varphi(\cdot,r)\theta(r)\dd r$ locally, and Cauchy--Schwarz.
\whyworks{convexity makes every fibre of the projection connected.}
\end{proof}
\uses{B22.}

%---------------------------------------------------------------------
\begin{dimcard}{B25}{C12-T1}{Factor decomposition on rectangles and convex sets}{\proved}{}
On a rectangle (and on every convex set), statements (i)--(iv) of A13 hold.
\end{dimcard}
\begin{proof}
\step{1}{Reduction}
Rectangles and convex sets are simply connected, so this is a special case of A13. The original proof replaces A12 by B23.
\whyworks{as for A13.}
\end{proof}
\uses{A13, B23.}

%---------------------------------------------------------------------
\begin{dimcard}{B26}{C12-P1}{Constructive projectors}{\proved}{}
(a) For $s=2$, $Q_2(\partial)u\in\ker Q_1$, and the components $u_1$ are exactly the solutions of $\{Q_1w=0,\ Q_2w=Q_2u\}$, unique up to $\ker Q_1\cap\ker Q_2$. (b) If all factor curves are disjoint, $u_j=E_j(\partial)u$ with Chinese remainder idempotents $E_j\in\ring$, and $u\in H^m\Rightarrow u_j\in H^{m-\deg E_j}$. (c) If $r=AQ_1+BQ_2\in\C[\xi]$ and $s=A'Q_1+B'Q_2\in\C[\tau]$, every component satisfies
\[
r(\partial_x)u_1=B(\partial)Q_2(\partial)u,\qquad s(\partial_t)u_1=B'(\partial)Q_2(\partial)u .
\]
\end{dimcard}
\begin{proof}
\step{1}{(a)}
$Q_1Q_2u=0$; if $u=u_1+u_2$ then $Q_2u_1=Q_2u$, and two solutions differ by an element of $\ker Q_1\cap\ker Q_2$.
\step{2}{(b)}
From $A_iQ_i+B_i\prod_{j\ne i}Q_j=1$ put $E_i=B_i\prod_{j\ne i}Q_j$; then $Q_iE_i=B_iP$ and $\sum_iE_i\equiv1\pmod P$. Each derivative costs one Sobolev order.
\step{3}{(c)}
$r(\partial_x)u_1=(AQ_1+BQ_2)u_1=BQ_2u_1=BQ_2u$.
\whyworks{every element of the ideal $(Q_1,Q_2)$ turns into an equation that a component must satisfy.}
\end{proof}
\uses{A11.}

%---------------------------------------------------------------------
\begin{dimcard}{B27}{C12-T2}{Vandermonde formula}{\proved}{}
For $P=\prod_{j=1}^k(\tau-a_j(\xi))$ with distinct $a_j$, $V(\xi)=[a_j^l]$, $\Delta=\det V=\prod_{i<j}(a_j-a_i)$:
\[
\Delta(\partial_x)\,u_j=\sum_{l=0}^{k-1}\mathrm{adj}(V)_{jl}(\partial_x)\,\partial_t^lu,\qquad \delta(P)=\deg\Delta .
\]
The split is unique iff $\Delta$ is a nonzero constant, and then $u_j=\Delta^{-1}\sum_l\mathrm{adj}(V)_{jl}(\partial_x)\partial_t^lu$.
\end{dimcard}
\begin{proof}
\step{1}{The jet in terms of the components}
$\partial_tu_j=a_j(\partial_x)u_j$, so $\partial_t^lu=\sum_ja_j(\partial_x)^lu_j$, that is $U=V(\partial_x)\mathbf u$.
\step{2}{Inverting}
Multiply by $\mathrm{adj}\,V$ and use $\mathrm{adj}(V)V=\Delta I$.
\step{3}{The ambiguity}
$\ring/(\tau-a_i,\tau-a_j)\cong\C[\xi]/(a_i-a_j)$, so $\delta(P)=\sum_{i<j}\deg(a_i-a_j)=\deg\Delta$.
\whyworks{the time derivatives of $u$ are the components weighted by a Vandermonde matrix in the symbols.}
\end{proof}
\uses{A13.}

\begin{example}{Wave}
For $(\tau-\xi)(\tau+\xi)$, $\Delta=-2\xi$ and the formula is d'Alembert's $\partial_xu_+=\tfrac12(\partial_xu+\partial_tu)$.
\end{example}

%---------------------------------------------------------------------
\begin{dimcard}{B28}{C12-P2}{Reducibility in second-order time; first order}{\proved{} (Lean-checked algebra)}{}
\phantomsection\label{dim:C12-P3}%
(C12-P2) $P=\tau^2+b(\xi)\tau+c(\xi)$ is reducible iff $D=b^2-4c$ is a square $d^2$ in $\C[\xi]$; then $P=(\tau-a_+)(\tau-a_-)$, $a_\pm=(-b\pm d)/2$, and $\delta(P)=\deg d$. (C12-P3) An evolution $\tau-a(\xi)$ has exactly one component.
\end{dimcard}
\begin{proof}
\step{1}{(P2)}
A factor of degree $0$ in $\tau$ divides the leading coefficient $1$, so a nontrivial factorisation has two monic linear factors $\tau-\alpha$, $\tau-\beta$ with $\alpha+\beta=-b$, $\alpha\beta=c$, $(\alpha-\beta)^2=D$. And $\delta=\deg(a_+-a_-)$ by B27.
\step{2}{(P3)}
A primitive polynomial of degree one in $\tau$ is irreducible (Gauss).
\whyworks{a quadratic in $\tau$ factors exactly when its discriminant is a square.}
\end{proof}
\uses{B27.}

\begin{example}{Klein--Gordon and telegraph}
For $\tau^2-\xi^2+m^2$, $D=4(\xi^2-m^2)$ is a square iff $m=0$; for $\tau^2+a\tau-\xi^2$, $D=a^2+4\xi^2$ is a square iff $a=0$. So both are single components.
\end{example}

%---------------------------------------------------------------------
\begin{dimcard}{B29}{C12-T4}{Interior regularity of the components}{\proved}{}
\phantomsection\label{dim:C12-L4}%
\textbf{(C12-L4)} If $p(\partial_x)w$ and $q(\partial_t)w$ lie in $H^k_{loc}$, with $\deg p=\deg q=D\ge1$, then $w\in H^{k+D}_{loc}$. \textbf{(C12-T4)} For coprime, non-comaximal $Q_1,Q_2$, with $r,s$ as in B26(c), $D=\deg r\deg s$ and $e$ the degree of the operators $BQ_2$, $B'Q_2$: every component of $u\in H^m_{loc}\cap\ker Q_1Q_2$ lies in $H^{m-e+D}_{loc}$.
\end{dimcard}
\begin{proof}
\step{1}{(L4) A Fourier estimate}
$|p(i\kappa)|\ge c|\kappa|^D-C(1+|\kappa|)^{D-1}$, and likewise for $q$, so
\[
(1+\omega^2+\kappa^2)^{D/2}\le C\big(1+|p(i\kappa)|+|q(i\omega)|\big).
\]
Localise with cutoffs; each commutator costs one order less, so finitely many steps reach $H^{k+D}_{loc}$.
\step{2}{(T4)}
Apply B26(c) with $p=r^{\deg s}$, $q=s^{\deg r}$ in the ideal, then (L4).
\whyworks{an ODE in $x$ and an ODE in $t$ together form a quasi-elliptic pair, which regularises.}
\end{proof}
\uses{B26.}

Interior regularity is never lost; A15 shows the loss happens only at the boundary.

%---------------------------------------------------------------------
\begin{dimcard}{B30}{C12-P4}{Two cases where the $H^m$ split holds}{\proved}{}
(a) $P=\tau\xi$ on a bounded rectangle, every $m$: $c_F=0$ after removing the constants. (b) The wave $\tau^2-\xi^2$ on the square $(0,1)^2$, every $m$: $\|u_+\|^2+\|u_-\|^2\le C\|u\|^2$ after normalisation, and $c_F<1$.
\end{dimcard}
\begin{proof}
\step{1}{(a)}
By B25, $u=A(t)\otimes1+1\otimes B(x)$. Integrating against $\theta(x)$ shows $A\in H^m$, hence $B\in H^m$. If $\int A=0$, the $H^m$ product of $A\otimes1$ and $1\otimes B$ reduces to the $\alpha=0$ term $\int A\int B=0$.
\step{2}{(b) Characteristic coordinates}
With $s=x+t$, $r=x-t$ the square becomes the rhombus $|r|<s<2-|r|$, and $u=F(s)+G(r)$ (B24).
\step{3}{(b) Five pieces}
On the central square $(\tfrac12,\tfrac32)\times(-\tfrac12,\tfrac12)$, averaging against test functions controls $F$ and $G$ by $u$. On each of the four corner triangles one of the two profiles is already controlled (its argument stays in the central range), so the other is $u$ minus it. The five pieces cover the rhombus.
\whyworks{when the two families of characteristics are transversal, every point sees one profile on a long chord.}
\end{proof}
\uses{B24, B25.}

%---------------------------------------------------------------------
\begin{dimcard}{B31}{C12-P5}{Variable coefficients: LCLM instead of the product}{\modulo{Ore 1933}}{}
For $L_1,L_2\in\C(x)[\partial]$ on an interval $J$ without singular points: $\ker\mathrm{LCLM}(L_1,L_2)=\ker L_1+\ker L_2$ and $\ker\mathrm{GCRD}(L_1,L_2)=\ker L_1\cap\ker L_2$. The product $L_1L_2$ does not work: $\ker((x\partial-1)\partial)\ne\ker(x\partial-1)+\ker\partial$.
\end{dimcard}
\begin{proof}
\step{1}{Inclusions}
$G=UL_1+VL_2$, $L_i=W_iG$, $M=AL_1=BL_2$ hold as identities of operators, and continuity extends them across the poles of the rational coefficients. So $\ker G=\ker L_1\cap\ker L_2$ and $\ker L_i\subseteq\ker M$.
\step{2}{Dimensions}
By Ore's identity $\mathrm{ord}M+\mathrm{ord}G=\mathrm{ord}L_1+\mathrm{ord}L_2$ \cf, so $\dim(\ker L_1+\ker L_2)=\mathrm{ord}M=\dim\ker M$.
\whyworks{in a noncommutative ring the least common left multiple plays the role of the product.}
\end{proof}
\uses{Ore 1933.}

%=====================================================================
\section{Regions with holes}

\begin{dimcard}{B32}{OPF-T4}{Cohomological criterion}{\pw{} ((i), (ii) need only sheaf theory)}{thm:cohom}
Let $\Omega$ be connected, $\mathcal N$ the $\delta(P)$-dimensional space of global exponential tuples with sum $0$, $\mathcal K_j$, $\mathcal K_P$ the solution sheaves. (i) $0\to\underline{\mathcal N}_\Omega\to\bigoplus_j\mathcal K_j\xrightarrow{\Sigma}\mathcal K_P\to0$ is exact. (ii) $\operatorname{coker}\Sigma_\Omega\cong\ker\big[\iota:H^1(\Omega;\C)\otimes\mathcal N\to\bigoplus_jH^1(\Omega;\mathcal K_j)\big]$; so the decomposition holds on $\Omega$ iff $\iota$ is injective, and $\dimc\operatorname{coker}\le\delta(P)\,b_1(\Omega)$. (iii) For the Laplacian on a disc with $k$ holes the cokernel has dimension exactly $k$ \cf.
\end{dimcard}
\begin{proof}
\step{1}{(i) Exactness on small discs}
On every disc A13 gives surjectivity and kernel $\mathcal N$; restriction between discs is the identity on $\mathcal N$, so the kernel sheaf is the constant sheaf $\underline{\mathcal N}$.
\step{2}{(ii) The long exact sequence}
\[
0\to\mathcal N\to\bigoplus_j\ker Q_j\xrightarrow{\Sigma_\Omega}\ker P\xrightarrow{\ \partial\ }H^1(\Omega;\underline{\mathcal N})\xrightarrow{\ \iota\ }\bigoplus_jH^1(\Omega;\mathcal K_j).
\]
So $\operatorname{coker}\Sigma_\Omega\cong\mathrm{im}\,\partial=\ker\iota$, and $H^1(\Omega;\underline{\mathcal N})\cong H^1(\Omega;\C)\otimes\mathcal N$ has dimension $\delta(P)\,b_1(\Omega)$.
\step{3}{(iii)}
Elliptic factors are surjective on every planar open set \cf, so $H^1(\Omega;\mathcal K_j)=0$ and the cokernel is all of $H^1\otimes\mathcal N$; for the Laplacian $\delta=1$, with basis $\log|z-a_l|^2$.
\whyworks{the obstruction to splitting is the monodromy of the shared exponentials around the holes.}
\end{proof}
\uses{A13; sheaf cohomology.}

%=====================================================================
\section{Closedness of sums of kernels}

\begin{dimcard}{B33}{P2-T3}{Two subspaces: exact constants}{\proved}{}
\phantomsection\label{dim:P2-C2}%
For closed $\mathcal H_1,\mathcal H_2$ with intersection $\mathcal I$: $\ker S=\{(z,-z):z\in\mathcal I\}$; on $(\ker S)^\perp$ the spectrum of $S^*S$ lies in $[1-c_F,1+c_F]\cup\{2\}$ and contains $1\pm c_F$. The minimal-energy split is stable iff $c_F<1$, with optimal constant $(1-c_F)^{-1}$. \textbf{(P2-C2)} For $N\ge2$ subspaces the same holds with $\gamma_N$; Tikhonov regularisation converges to the minimal-energy split, and its norm blows up when $\gamma_N=0$.
\end{dimcard}
\begin{proof}
\step{1}{Spectrum}
This is A17, Steps~2--3: $S^*S=I+K$ on $A\oplus B$ with $\|K\|=c_F$ and symmetric spectrum, and $S^*S=2$ on the diagonal of $\mathcal I$.
\step{2}{Tikhonov}
$\omega_\alpha=(S^*S+\alpha)^{-1}S^*u$; in the spectral resolution of $S^*S$ it converges to $S^\dagger u$ when $u\in\mathrm{im}\,S$, and $\|\omega_\alpha\|\to\infty$ for $u\in\overline{\mathrm{im}\,S}\setminus\mathrm{im}\,S$.
\whyworks{the sum map is a block operator whose off-diagonal part has norm equal to the cosine of the angle.}
\end{proof}
\uses{A17.}

%---------------------------------------------------------------------
\begin{dimcard}{B34}{P2-C1}{Where regularising the factor split fails}{\proved}{}
(a) For two parallel transports on a convex $\Omega$ in $L^2$ the sum is dense in $\ker fg$; if the chords shrink, some $u$ has no $L^2$ components, and for every such $u$ the Tikhonov split diverges while its residual tends to $0$. (b) On the square, for the function of A15, either every regularisation leaves a residual, or $c_F=1$ and the regularised components diverge.
\end{dimcard}
\begin{proof}
\step{1}{(a) Density}
In coordinates $(t,s)$, $u=e^{d_1t}F(s)+e^{d_2t}G(s)$; truncating $F$, $G$ to compact intervals of $s$ gives elements of $\mathcal H_1+\mathcal H_2$ converging to $u$.
\step{2}{Divergence}
B33 with $c_F=1$ (A16).
\whyworks{dense but not closed means: approximable, but only by pieces of unbounded size.}
\end{proof}
\uses{A16, B33.}

%---------------------------------------------------------------------
\begin{dimcard}{B35}{OPF-T5}{Density via contracting symmetries}{\pw}{}
Let $\Omega$ be bounded and convex, and suppose affine maps $\Phi_\lambda(y)=o+A_\lambda(y-o)$, $A_\lambda\to I$, satisfy $P(A_\lambda^{\mathsf T}z)=\kappa_\lambda P(z)$ and $\Phi_\lambda(\bar\Omega)\subset\Omega$. Then $\sum_j\ker_{H^m}Q_j$ is dense in $\ker_{H^m}P$ for every $m$. This covers homogeneous $P$ (the wave on every convex set) and quasi-homogeneous $P$ on axis-parallel rectangles (the beam).
\end{dimcard}
\begin{proof}
\step{1}{Shrink, mollify, decompose}
$u\circ\Phi_\lambda\to u$ in $H^m$ and solves $P$ (the symmetry). Mollify on a convex neighbourhood of $\Phi_\lambda(\bar\Omega)$, pull back: a smooth solution on a convex neighbourhood of $\bar\Omega$, which splits by A13 into components smooth up to the boundary.
\whyworks{a symmetry of $P$ pulls a solution slightly inward without leaving the solution space.}
\end{proof}
\uses{A13.}

This is superseded on convex sets by A18, which needs no symmetry.

%---------------------------------------------------------------------
\begin{dimcard}{B36}{OPF-P6}{Parallel families: pairwise closed iff totally closed}{\pw}{}
For $Q_j=\partial_t+c\partial_x-d_j$ ($d_j$ distinct) on a bounded convex $\Omega$ in $L^2$, equivalent: (a) the chords have $\inf\ell>0$; (b) $\sum_j\mathcal H_j$ closed; (c) every pairwise sum closed; (d) some pairwise sum closed.
\end{dimcard}
\begin{proof}
\step{1}{Structure}
The factors are comaximal, so the split is unique, $u=\sum_je^{d_jt}F_j(s)$.
\step{2}{(a)$\Rightarrow$(b)}
The $N\times N$ Gram matrix of $\{e^{d_jt}\}$ on chords of length $\ge\ell_0$ is uniformly positive definite, as in A16, Step~2.
\step{3}{(d)$\Rightarrow$(a)}
If $\inf\ell=0$, A16, Step~4, applied to any pair shows that pair is not closed.
\whyworks{for parallel transports everything happens chord by chord, and all pairs see the same chords.}
\end{proof}
\uses{A16.}

%---------------------------------------------------------------------
\begin{dimcard}{B37}{OPF-P7}{Separable families: the sum is always closed}{\pw}{}
If every factor depends on $t$ alone or on $x$ alone, $P=A(\tau)B(\xi)$ on a bounded rectangle, then $\sum_j\ker_{H^m}Q_j$ is closed in $H^m$ for every grouping.
\end{dimcard}
\begin{proof}
\step{1}{A bounded projector}
With a basis $e_l$ of $\ker A(\partial_t)$ and biorthogonal $\psi_l\in C_c^\infty(I)$, $\pi_Au=\sum_le_l(t)\int u(t',x)\psi_l(t')\dd t'$ is a bounded idempotent on $H^m$ with image $\ker_{H^m}A(\partial_t)$; likewise $\pi_B$.
\step{2}{Closed sum}
The $t$-factors sum to $\mathrm{im}\,\pi_A$ and the $x$-factors to $\mathrm{im}\,\pi_B$; two ranges of commuting bounded projectors have a closed sum, $\mathrm{im}(\pi_A+\pi_B-\pi_A\pi_B)$.
\whyworks{separated variables give commuting projectors.}
\end{proof}
\uses{B23.}

%---------------------------------------------------------------------
\begin{dimcard}{B38}{OPR2-T2}{Real transports in $L^2$}{\pw{} \emph{(condensed)}}{prop:transports}
For real transports $Q_j=v_j\cdot\nabla-d_j$ on a bounded convex $\Omega$, grouped in classes of parallel directions, assume that of any two classes at least one has chords bounded below. Then $\sum_j\ker_{L^2}Q_j$ is closed iff every class with at least two factors has chords bounded below; in particular pairwise closed implies totally closed.
\end{dimcard}
\begin{proof}
\step{1}{Class projectors}
For a class with chords bounded below, B36 gives a closed sum with a chord-by-chord projector whose kernel is bounded.
\step{2}{Compact cross products}
For two transversal classes, the product of their projectors has a kernel integrable over the transversal chords, hence is compact, under the hypothesis on the chords.
\step{3}{Closed sum}
A finite sum of closed subspaces whose projectors have compact pairwise products and trivial intersections is closed (Fredholm argument). If a class with two factors has shrinking chords, B36 shows the class sum itself is not closed.
\whyworks{transversal transports interact compactly, so only parallel ones can destroy closedness.}
\end{proof}
\uses{B36; research file \texttt{work/OP-R2}.}

%---------------------------------------------------------------------
\begin{dimcard}{B39}{OPR2-T1}{Heat times transport: density}{\pw}{}
For $P=(\tau-\xi^2)(\tau+c\xi)$ on a rectangle, $\ker_{H^m}(\partial_t-\partial_x^2)+\ker_{H^m}(\partial_t+c\partial_x)$ is dense in $\ker_{H^m}P$ for every $m$ and $c$; so the components lie in $H^m$ iff $c_F<1$.
\end{dimcard}
\begin{proof}
\step{1}{From A18}
A rectangle is bounded and convex, so A18 applies directly. (The original proof used local inward translations at the corners and the hypoellipticity of the heat operator.)
\whyworks{as for A18.}
\end{proof}
\uses{A18.}

%---------------------------------------------------------------------
\begin{dimcard}{B40}{OPF5B-P8}{The tangent model of a wedge is closed}{\pw}{}
On the infinite wedge $C=\{r\omega(\varphi):\ 0<\varphi<\beta\}$, $\beta<\pi$, with pairwise non-parallel homogeneous transports whose lines all cut $C$ in bounded chords, there is $\gamma_*>0$ with $\|\sum_jf_j\|^2\ge\gamma_*\sum_j\|f_j\|^2$ for $f_j\in\mathcal H_j$.
\end{dimcard}
\begin{proof}
\step{1}{Structure}
$f\in\mathcal H_j$ iff $f(y)=F(n_j\cdot y)$ with $F\in L^2(s\dd s)$, since the chords have length proportional to $s$.
\step{2}{Mellin transform}
The Mellin transform in $r$ is unitary and maps $\mathcal H_j$ to $\{F^\sharp(\sigma)\,a_j(\varphi)^{-1+i\sigma}\}$, with $\|a_j^{-1+i\sigma}\|_{L^2(0,\beta)}$ independent of $\sigma$.
\step{3}{Fibrewise Gram matrices}
The norm of a sum becomes an integral over $\sigma$ of quadratic forms with the Gram matrix $G(\sigma)$ of the functions $a_j^{-1+i\sigma}$, which is positive definite and has a positive infimum of its smallest eigenvalue (the limits $\sigma\to\pm\infty$ decouple by oscillation).
\whyworks{dilation invariance diagonalises the problem, leaving one finite Gram matrix per frequency.}
\end{proof}
\uses{none.}

%=====================================================================
\section{Complements: variable-coefficient identification}

\begin{dimcard}{B41}{OPF5B-T2}{Saturated criterion for non-monic $P$}{\pw}{}
\phantomsection\label{dim:OPF5B-P3}\phantomsection\label{dim:OPF5B-C4}%
Let $P=\sum_{j\le m}p_j\partial_t^j$ with $p_m\in\mathfrak A\setminus0$, $\mathfrak K$ the Ore localisation. For $Pu=0$, equivalent: (a$'$) $\Ann_{\mathfrak D}(u)\subseteq\mathfrak KP$; (b) the $t$-jet is free over $\mathfrak A$; (c) $\mathrm{rank}_{\mathfrak A}\mathfrak Du=m$; (d$'$) no $Q\notin\mathfrak KP$ has $Qu$ torsion. \textbf{(OPF5B-P3)} If $p_m=\pi(x)$ is a function and $\mathfrak D/\mathfrak DP$ has no $\pi$-torsion, B20 holds verbatim. \textbf{(OPF5B-C4)} The radial heat kernel $t^{-1}e^{-x^2/4t}$ identifies the radial heat operator $x\partial_t-\partial_xx\partial_x$; the Bessel modes $e^{-k^2t}J_0(kx)$ do not.
\end{dimcard}
\begin{proof}
\step{1}{Localise}
In $\mathfrak K$, $P=p_m\hat P$ with $\hat P$ monic; divide by $\hat P$ and clear denominators with the Ore condition to turn a non-multiple annihilator into a nontrivial relation of degree $<m$, as in B20.
\step{2}{Rank}
Localisation is exact: $\mathrm{rank}\,\mathfrak Du=\dim_K\mathfrak K/\mathfrak K\Ann(u)\le m$, with equality iff $\Ann(u)\subseteq\mathfrak KP$.
\step{3}{(C4)}
The right action of $P_r$ on $\mathfrak D/x\mathfrak D\cong\C[\eta,T]$ is $f\mapsto Tf'-\eta f-\eta^2f'$, injective (leading coefficient $-(d+1)c_d$), so there is no $x$-torsion. The heat kernel is not killed by any nonzero operator in $x$ alone (its derivatives are polynomials of growing degree times the Gaussian); the Bessel modes satisfy Bessel's ODE.
\whyworks{a non-monic operator behaves like a monic one after inverting its leading coefficient, provided nothing is lost to torsion.}
\end{proof}
\uses{B20.}

%---------------------------------------------------------------------
\begin{dimcard}{B42}{OPF5B-T6}{Repeated socle: no finite test}{\pw}{}
\phantomsection\label{dim:OPF5B-L5}\phantomsection\label{dim:OPF5B-P7}%
Assume some solution $u^*$ identifies $P$ (monic) and $\mathrm{End}_{\mathfrak D}(\mathfrak D/\mathfrak DP)$ contains infinitely many non-injective endomorphisms $E_\lambda$ with pairwise trivially intersecting kernels. Then for every finite set of tests $\mathcal T_f$ the criterion ``$\Ann(u)=\mathfrak DP$ iff $Tu$ is not torsion for all $T\in\mathcal T_f$'' fails for some solution. The hypothesis holds for $P=P_1P_2$, $P_2=P_1+c$, whenever a nonzero $J$ with $P_1J\in\mathfrak DP_2$ exists (example: $\theta=x\partial_x$, $J=\partial_x$).
\end{dimcard}
\begin{proof}
\step{1}{Solutions from endomorphisms (OPF5B-L5)}
$E(\bar Q)=\overline{QR}$ with $PR\in\mathfrak DP$; $u_R=Ru^*$ solves $P$, and $u_R$ identifies iff $E_R$ is injective.
\step{2}{Beating a finite test}
Each test $T$ lies in at most one kernel $\ker E_\lambda$; choose $\lambda$ outside the finitely many used. Then $u_\lambda$ passes every test but does not identify.
\step{3}{The example (OPF5B-P7)}
$\mathfrak D/\mathfrak DP\cong\mathfrak D/\mathfrak DP_1\oplus\mathfrak D/\mathfrak DP_2$ via $\bar Q\mapsto(\bar Q,\bar Q)$, and $E_\lambda(x_1,x_2)=(0,x_2+\lambda Jx_1)$ has kernel $\{(x_1,-\lambda Jx_1)\}$, pairwise trivially intersecting.
\whyworks{a repeated simple factor offers infinitely many ways to fail, more than any finite list of tests can detect.}
\end{proof}
\uses{B20.}

\section{Results used without proof}

\listcard{C12-P6}{No orthogonality of components at fixed time in general.}{\cf}{Shohat, Amer.\ J.\ Math.\ (checked, CIT-AUDIT-2).}
\listcard{C12-Q2}{Closedness of the sum for $N\ge3$ general differential factors.}{\openq}{Solved for real transports (B38), separable families (B37), wedges in the tangent model (B40).}
\listcard{OPF-Q3}{Density for quasi-homogeneous $P$ on non-rectangular convex sets.}{closed by A18}{A18 covers every bounded convex set.}
